\documentclass[9pt,letterpaper]{article}
\usepackage[utf8]{inputenc}
\usepackage[spanish]{babel} 
\decimalpoint

%----- Configuración del estilo del documento------%
\usepackage{epsfig,graphicx}
\usepackage[left=2cm,right=2cm,top=1.8cm,bottom=2.3cm]{geometry}
\usepackage{fancyhdr}
\pagestyle{fancy}
\fancyhf{}
\fancyfoot{}
\graphicspath{{Figuras/}}
\usepackage{hyperref}
\hypersetup{colorlinks=true,linkcolor=blue,citecolor=blue,filecolor=blue,urlcolor=magenta,}

\fancyfoot[RO,LE]{\thepage} % Custom footer text
\fancyheadoffset[RO,LE]{0.01\textwidth}

%------ Paquetes matemáticos básicos --------%
\usepackage{amsmath}
\usepackage{amssymb}
\usepackage{amsthm}

%------ Texto aleatorio ----- %

\usepackage{lipsum}



\begin{document}
	
	%------ Encabezado -------- %
	
	\begin{center}
		\begin{minipage}{3cm}
			\begin{center}
				\includegraphics[height=3.4cm]{../Figuras/Logo_UNAM (1)}
			\end{center}
		\end{minipage}\hfill
		\begin{minipage}{10cm}
			\begin{center}
				{\scshape\LARGE \textbf{Universidad Nacional Autónoma de México} \par}
				{\scshape\Large Facultad de Ciencias\par}
				{\scshape\Large Coordinación de Física Biomédica\par}
			\end{center}
		\end{minipage}\hfill
		\begin{minipage}{3cm}
			\begin{center}
				\includegraphics[height=3.4cm]{../Figuras/Logo_FC (1)}
			\end{center}
		\end{minipage}
	\end{center}
	
	\rule{17cm}{0.1mm}
	
	%------ Fin de encabezado -------- %
	\hspace{0.5cm}
	
	\parbox{\textwidth}{\raggedleft Ciudad Universitaria,  de  de 2025.}
	
	\hspace{1cm}
	
	\vspace{0.5cm}
	
	Comisión de Servicio Social
	
	Coordinación de Física Biomédica
	
	Facultad de Ciencias
	
	Universidad Nacional Autónoma de México
	
	P R E S E N T E
	
	\vspace{0.5cm}
	
	
	Por medio de la presente, expongo las razones por las que me retrasé en la entrega de documentos para finalización del
	Servicio Social, que realicé dentro del programa Apoyo a la Investigación, con el título ``Cálculo de las secciones transversales de extinción, absorción y esparcimiento de elipsoides en la aproximación cuasi-estática como primera aproximación de eritrocitos'' con clave 2024-12/12-837. El proyecto lo realicé bajo la supervisión del Dr. Alejandro Reyes Coronado, en el Departamento de Física, cubículo 407 en la Facultad de Ciencias de la UNAM, el cual estaba planeado para llevarse a cabo durante el periodo comprendido entre el 23 de febrero de 2024 al 23 de septiembre del 2024. \\
	
	
	\noindent Las principales razones fueron asuntos familiares y académicos. 
	
	\vspace{1cm}
	\noindent Agradeciendo su comprensión, sin más, me despido con un cordial saludo.
	
	{\vspace{2.55cm}\begin{tabular} { c}
			\setlength{\tabcolsep}{15pt}
			\renewcommand{\arraystretch}{1}
			\noindent\rule{5.5cm}{0.4pt}\qquad \\
			
			\qquad  \textbf{Dana Larissa Luna González} \qquad \\
			\qquad Estudiante de Física Biomédica  \qquad \\ \qquad 
			No. de cuenta: 421122680\qquad \\  
			\qquad  Tel.: 776 101 4262 \qquad \\
			\qquad dana.larissalg@ciencias.unam.mx \qquad \\
			
		\end{tabular}
	}
	
	{\vspace{-2.53cm}\hspace{7cm}\begin{tabular} { c}
			\setlength{\tabcolsep}{15pt}
			\renewcommand{\arraystretch}{1}
			\noindent\rule{5.5cm}{0.4pt}\qquad \\
			
			\qquad  \textbf{Alejandro Reyes Coronado} \qquad \\
			\qquad Profesor Titular C de tiempo completo  \qquad \\  
			\qquad Departamento de Física, Facultad de Ciencias, UNAM\qquad \\ 
			\qquad  Tel.: (55) 5622 4968 \qquad \\
			\qquad coronado@ciencias.unam.mx \qquad \\
			
		\end{tabular}
		
	}
	
	
	
	
	
	
	
	
	
\end{document}


